\begin{exercise}{46}
Let $Y\to X$ be an inclusion of topological spaces. When is the induced map
$\mathcal S(Y)\to\mathcal S(X)$ a subcategory? When is it a full subcategory?

\begin{proof}
The answer to both questions is: exactly when $Y$ is a union
of path components of $X$.
Both singular complexes are Kan by Lemma~1.1.47.
Identify $\mathcal S(Y)$ with the singular simplices of $X$
whose images lie in $Y$; corestriction to the subspace $Y$
is continuous. If $Y$ is such a union, any singular simplex
with a vertex in $Y$ lies entirely in $Y$, since its domain
is path-connected. Thus $\mathcal S(Y)$ consists precisely
of the simplices of $\mathcal S(X)$ whose vertices lie in $Y$.
It is the full subcategory on $Y$ by Definition~1.2.76.

Conversely, suppose $\mathcal S(Y)$ is a subcategory and let
$p\colon[0,1]\to X$ satisfy $p(0)\in Y$.
Compose $p$ with the map
$\sigma\colon\Delta^2_{\mathrm{top}}\to[0,1]$
from Exercise~\ref{ex:45}.
The resulting singular $2$-simplex has both spine edges
constant at $p(0)$, hence in $\mathcal S(Y)$.
Definition~1.2.76 forces the whole simplex to lie in
$\mathcal S(Y)$; evaluating at the barycenter gives
$p(1)\in Y$. Hence every path component meeting $Y$ is
contained in $Y$.
A full subcategory is in particular a subcategory, so this
also proves necessity for the second question.
\end{proof}
\end{exercise}
